\apendice{Anexo de sostenibilización curricular}

\section{Introducción}
El presente anexo recoge una reflexión crítica sobre la alineación del
\ACRshort{tfg} con los principios de sostenibilidad y responsabilidad social,
en concordancia con las directrices establecidas por la \acrfull{crue}.

La ingeniería del software, tradicionalmente percibida como una disciplina
aséptica, tiene un impacto profundo y directo en el consumo de recursos
naturales, el desarrollo socioeconómico y la equidad social. Por tanto, el
diseño, arquitectura y desarrollo de este sistema de análisis temático se ha
llevado a cabo bajo el prisma de sostenibilidad ambiental, social y económica.

A continuación, en la Tabla~\ref{tabla:ods_resumen}, se resumen las principales
decisiones adoptadas y su impacto directo en la sostenibilidad.

\tablaSmall{Resumen del impacto de las decisiones técnicas en la sostenibilidad.}
{p{0.7\textwidth} p{0.3\textwidth}}
{ods_resumen}
{
    \textbf{Decisión del proyecto}      & \textbf{Dimensión Sostenible} \\
}
{
    Uso de IA ligera local y Parquet    & Ambiental        \\
    Desarrollo de herramienta analítica & Social           \\
    Anonimización de datos (RGPD)       & Social / Ética   \\
    Licenciamiento de código abierto    & Económica        \\
}

\section{Sostenibilidad Ambiental}
El auge reciente de la Inteligencia Artificial y de los modelos de
Procesamiento de Lenguaje Natural (\acrshort{nlp}) ha traído consigo un
incremento alarmante en la huella de carbono, debido a la ingente cantidad de
recursos computacionales y energéticos necesarios para su entrenamiento. En
este proyecto, se han aplicado principios de \textit{Green AI} (IA verde) para
minimizar drásticamente el impacto medioambiental.

En primer lugar, se ha evitado la dependencia de modelos de lenguaje cerrados
en la nube que requieren centros de datos masivos. En su lugar, se ha optado
por modelos probabilísticos ligeros (como \acrshort{lda}) y modelos basados en
Transformers de tamaño reducido (como BERTopic y Top2Vec) que pueden ser
ejecutados y afinados localmente utilizando un hardware doméstico estándar.

Además, la arquitectura de persistencia de datos ha sido diseñada para ser
altamente eficiente. La transición de formatos tradicionales en texto plano
(como JSON o \acrshort{csv}) hacia el uso exclusivo de Apache Parquet ha
permitido comprimir drásticamente el espacio de almacenamiento. Parquet, al ser
un formato columnar optimizado, minimiza las operaciones de lectura y escritura
en disco (I/O), reduciendo notablemente el consumo energético del procesador y
el disco durante las pesadas fases de tratamiento de texto.

\section{Sostenibilidad Social y Ética}
Desde el punto de vista del impacto humano, el proyecto contribuye de forma
directa al Educación de Calidad, Industria, Innovación e Infraestructura. La
herramienta desarrollada tiene un fin puramente académico y divulgativo:
proporcionar a investigadores, docentes y futuros estudiantes una plataforma
analítica intuitiva para comprender las tendencias, tecnologías y temáticas que
se están investigando en el ámbito universitario. Al extraer el conocimiento
latente de memorias de \Acrshort{tfg}, se facilita la transferencia de
conocimiento, se inspiran nuevas líneas de investigación para el alumnado y se
evita la redundancia académica.

\section{Sostenibilidad Económica}
El ecosistema de este desarrollo se ha sustentado de forma exclusiva sobre
tecnologías libres y gratuitas. El \textit{backend} (Python, Pandas, Optuna) y
el motor de visualización (Streamlit) no exigen licencias comerciales de pago.

Asimismo, la ligereza de la interfaz web permite que el producto final pueda
ser albergado en plataformas en la nube de capa gratuita (como \textit{GitHub
    Container Registry} o \textit{Streamlit Community Cloud}), garantizando una
disponibilidad analítica universal a coste cero de mantenimiento operativo.

\section{Conclusión}
En conclusión, la elaboración de este \Acrshort{tfg} ha supuesto una
oportunidad invaluable para interiorizar y poner en práctica las competencias
transversales de sostenibilidad impulsadas por la \acrshort{crue}. Las
decisiones de diseño técnico adoptadas desde la selección del formato Parquet
para reducir el consumo energético, hasta la adopción de una licencia de código
abierto para democratizar su uso y la anonimización proactiva de los datos
recogidos demuestran que es posible concebir productos de software que sean
técnicamente robustos a la par que social, económica y medioambientalmente
responsables. Este proyecto contribuye a través de una praxis ética y
respetuosa con su entorno material y humano.
